\documentclass[UTF8,a4paper,12pt]{ctexart}

\usepackage[margin=2.5cm]{geometry}
\usepackage{amsmath,amssymb}
\usepackage{booktabs}
\usepackage{longtable}
\usepackage{array}
\usepackage{hyperref}
\usepackage{url}
\usepackage{xcolor}

\hypersetup{
  colorlinks=true,
  linkcolor=black,
  urlcolor=blue
}

\setlength{\parindent}{2em}
\setlength{\parskip}{0.35em}
\renewcommand{\arraystretch}{1.18}

\title{问题三：基于清水池停留时间分布与数据驱动残差修正的出厂水浊度动态预测}
\author{}
\date{}

\begin{document}
\maketitle

\section{问题目标与数据特点}

问题三要求预测未来 1--12 小时的出厂水浊度（NTU），并给出 2026 年 2 月 1 日、2 月 10 日、2 月 20 日 7 点至 19 点的预测结果。由于原水从取水、混凝、沉淀、过滤到清水池和出厂存在明显滞后，直接用当前变量同步拟合出厂 NTU 难以反映水力停留和混合过程。因此，本问采用“清水池水力停留时间分布（RTD）+ XGBoost 残差模型 + 自回归递推填补”的混合动态模型。

原始监测数据为每 2 小时一次，因此可直接由数据监督训练的预测步长为 2、4、6、8、10、12 小时。对于题目中“1--12 小时”的表述，本文在 2 小时网格上建立直接多步预测模型，并对 7 点至 19 点逐小时结果采用相邻 2 小时预测的线性插值补齐。

\section{建模流程}

整体流程如下：

\begin{enumerate}
  \item 读取 2025 年训练数据和 2026 年 1--3 月验证数据，统一时间索引；
  \item 构造清水池停留时间分布（RTD）加权特征，描述历史水质对当前出厂水状态的混合贡献；
  \item 构造滞后、差分、滚动统计和时间周期特征；
  \item 训练 2、4、6、8、10、12 小时六个预测步长的 XGBoost 残差模型；
  \item 用 2 小时模型按时间顺序递推填补 2026 年 2 月缺失的 NTU 状态；
  \item 在填补后的完整状态表上重新训练最终模型，并输出指定日期的预测结果；
  \item 在 2026 年 1 月和 3 月有真实 NTU 标签的样本上验证模型精度。
\end{enumerate}

\section{RTD 机理层}

清水池具有混合和滞后作用。模型首先根据清水池水位和原水流量构造等效平均停留时间：

\begin{equation}
\tau(t)=6 \times \frac{L(t)}{\bar L} \times \frac{\bar Q}{Q(t)},
\end{equation}

其中，$L(t)$ 为清水池水位 \texttt{C/W WELL LEVEL}，$Q(t)$ 为原水流量 \texttt{R/W FLOW}，$\bar L$ 与 $\bar Q$ 分别为对应变量的样本中位数。为避免极端流量导致不合理停留时间，$\tau(t)$ 被限制在 2--12 小时之间。

对任一输入变量 $X$，用 Gamma 型离散权重表示过去 24 小时对当前状态的贡献：

\begin{equation}
RTD_X(t)=\sum_{k=1}^{12} w_k(t)X(t-2k),
\end{equation}

其中 $w_k(t)\ge 0$ 且 $\sum_{k=1}^{12}w_k(t)=1$。权重由当前 $\tau(t)$ 决定，流量较大时权重更集中于近时刻，流量较小时权重向更早历史展开。

本模型对以下变量构造 RTD 特征：

\begin{itemize}
  \item 原水浊度：\texttt{R/W NTU}
  \item 滤后水浊度：\texttt{FILT. NTU}
  \item 矾投加量：\texttt{ALUM}
  \item 出厂水浊度状态：\texttt{NTU}
\end{itemize}

\section{物理基准项与残差模型}

基于质量守恒和清水池混合思想，构造物理基准项：

\begin{equation}
\widehat{NTU}_{phys}(t)
=0.45RTD_{FILT}(t)
+0.35RTD_{NTU}(t)
+0.15NTU(t-2h)
+0.05FILT(t).
\end{equation}

该基准项提供一个可解释的动态平滑预测：滤后水浊度反映进入清水池前的过程水质，历史出厂水 NTU 反映清水池内已有水体状态，当前滤后水浊度补充最新过程扰动。

由于混凝、过滤、加药和运行调控具有非线性，物理基准项无法完全描述所有误差。因此对每个预测步长 $h$ 训练一个 XGBoost 残差模型：

\begin{equation}
\widehat{NTU}(t+h)=\widehat{NTU}_{phys}(t)+g_h(\mathbf{x}_t),
\quad h\in\{2,4,6,8,10,12\}.
\end{equation}

其中 $\mathbf{x}_t$ 包括当前变量、RTD 特征、滞后特征、差分特征、滚动统计和时间周期特征。

\section{特征构造}

模型输入主要包括：

\begin{itemize}
  \item 当前工况：\texttt{RIVER LEVEL}、\texttt{R/W FLOW}、\texttt{R/W NTU}、\texttt{R/W CLR}、\texttt{R/W PH}、\texttt{FILT. NTU}、\texttt{C/W WELL LEVEL}、\texttt{ALUM}；
  \item RTD 特征：\texttt{RTD\_R\_W\_NTU}、\texttt{RTD\_FILT\_NTU}、\texttt{RTD\_ALUM}、\texttt{RTD\_NTU}；
  \item 滞后特征：2、4、6、8、10、12、18、24 小时滞后；
  \item 短期变化：2 小时差分；
  \item 局部波动：过去 12 小时滚动均值和标准差；
  \item 时间周期：小时和月份的正余弦编码。
\end{itemize}

\section{2 月 NTU 缺失的递推填补}

2026 年 2 月出厂 NTU 真实值缺失，而自回归模型需要历史 NTU 状态。若直接使用历史 NTU 滞后项，则 2 月首个缺失点后特征链会断裂。因此本文采用两阶段递推策略。

首先用 2025 年训练集训练初始 2 小时模型，然后从 2026 年 2 月 1 日 01:00 开始按时间顺序递推：

\begin{equation}
\widehat{NTU}(t+2h)=\widehat{NTU}_{phys}(t)+g_{2h}(\mathbf{x}_t).
\end{equation}

每得到一个预测值，就将其作为后续时刻的历史 NTU 状态继续构造滞后和 RTD 特征。递推填补覆盖 2026 年 2 月 1 日 01:00 至 2026 年 2 月 28 日 23:00，共 336 条记录，填补 NTU 范围为 0.141--0.806。

完成递推填补后，再基于完整状态表重新训练 2--12 小时最终模型，并输出题目所需日期的预测结果。

\section{模型验证}

由于 2026 年 2 月真实 NTU 缺失，验证指标仅在 2026 年 1 月和 3 月有标签样本上计算。整体验证结果如下：

\begin{table}[h]
\centering
\begin{tabular}{rrrr}
\toprule
预测步长 & MAE & RMSE & $R^2$ \\
\midrule
2h  & 0.0480 & 0.1171 & 0.7555 \\
4h  & 0.0630 & 0.1448 & 0.6264 \\
6h  & 0.0736 & 0.1657 & 0.5109 \\
8h  & 0.0798 & 0.1762 & 0.4468 \\
10h & 0.0788 & 0.1693 & 0.4891 \\
12h & 0.0803 & 0.1656 & 0.5111 \\
\bottomrule
\end{tabular}
\caption{2026 年 1 月和 3 月合并验证结果}
\end{table}

分月验证结果显示，2 小时预测在 1 月和 3 月均有较好表现；随着预测步长增加，1 月长步长预测精度下降较明显，而 3 月在 6--12 小时步长上保持相对稳定。这说明出厂水 NTU 的短期自回归状态对 2--4 小时预测最有效，长步长预测则更多依赖工艺过程变量和 RTD 特征。

\begin{longtable}{rrrrr}
\caption{按目标月份拆分的验证结果}\\
\toprule
步长 & 目标月份 & 样本数 & RMSE & $R^2$ \\
\midrule
\endfirsthead
\toprule
步长 & 目标月份 & 样本数 & RMSE & $R^2$ \\
\midrule
\endhead
2h  & 2026-01 & 371 & 0.0993 & 0.7699 \\
2h  & 2026-03 & 372 & 0.1326 & 0.4917 \\
4h  & 2026-01 & 370 & 0.1528 & 0.4562 \\
4h  & 2026-03 & 372 & 0.1364 & 0.4622 \\
6h  & 2026-01 & 369 & 0.1928 & 0.1360 \\
6h  & 2026-03 & 372 & 0.1334 & 0.4856 \\
8h  & 2026-01 & 368 & 0.2054 & 0.0214 \\
8h  & 2026-03 & 372 & 0.1414 & 0.4215 \\
10h & 2026-01 & 367 & 0.2026 & 0.0500 \\
10h & 2026-03 & 372 & 0.1283 & 0.5240 \\
12h & 2026-01 & 366 & 0.1915 & 0.1533 \\
12h & 2026-03 & 372 & 0.1355 & 0.4690 \\
\bottomrule
\end{longtable}

\section{指定日期预测结果}

以 7:00 为起报时刻，未来 2--12 小时预测结果如下：

\begin{table}[h]
\centering
\begin{tabular}{rrrr}
\toprule
日期 & 目标时刻 & 步长 & 预测 NTU \\
\midrule
2026-02-01 & 09:00 & 2h  & 0.2952 \\
2026-02-01 & 11:00 & 4h  & 0.3103 \\
2026-02-01 & 13:00 & 6h  & 0.3276 \\
2026-02-01 & 15:00 & 8h  & 0.3256 \\
2026-02-01 & 17:00 & 10h & 0.3519 \\
2026-02-01 & 19:00 & 12h & 0.3195 \\
2026-02-10 & 09:00 & 2h  & 0.1944 \\
2026-02-10 & 11:00 & 4h  & 0.1943 \\
2026-02-10 & 13:00 & 6h  & 0.1952 \\
2026-02-10 & 15:00 & 8h  & 0.1900 \\
2026-02-10 & 17:00 & 10h & 0.1811 \\
2026-02-10 & 19:00 & 12h & 0.1765 \\
2026-02-20 & 09:00 & 2h  & 0.1606 \\
2026-02-20 & 11:00 & 4h  & 0.1620 \\
2026-02-20 & 13:00 & 6h  & 0.1730 \\
2026-02-20 & 15:00 & 8h  & 0.1779 \\
2026-02-20 & 17:00 & 10h & 0.1683 \\
2026-02-20 & 19:00 & 12h & 0.1621 \\
\bottomrule
\end{tabular}
\caption{2026 年 2 月指定日期 7 点起报预测}
\end{table}

完整结果已保存至 \path{Question3/outputs/hybrid_dynamic_model/question3_answer.xlsx}。其中 \texttt{rolling\_7\_19} 为 2 小时滚动预测，\texttt{forecast\_from\_07} 为 7 点起报的 2--12 小时预测，\texttt{hourly\_interpolated} 为 7 点至 19 点逐小时插值结果。

\section{敏感性分析}

为评估原水水质突变和矾投加调整对预测结果的影响，本文在三个目标日期 7:00 起报样本上设置局部扰动场景：

\begin{itemize}
  \item 原水浊度增加 20\%；
  \item 原水浊度增加 50\%；
  \item 矾投加量增加 10\%；
  \item 矾投加量减少 10\%。
\end{itemize}

扰动不仅作用于当前变量，也同步作用于对应滞后特征和 RTD 加权特征。各场景相对于基准预测的变化统计如下：

\begin{table}[h]
\centering
\begin{tabular}{lrrr}
\toprule
扰动场景 & 最小变化 & 最大变化 & 平均变化 \\
\midrule
原水 NTU 增加 20\% & 0.0000 & 0.0009 & 0.0001 \\
原水 NTU 增加 50\% & 0.0000 & 0.1018 & 0.0061 \\
矾投加增加 10\% & 0.0000 & 0.0000 & 0.0000 \\
矾投加减少 10\% & 0.0000 & 0.0000 & 0.0000 \\
\bottomrule
\end{tabular}
\caption{指定日期局部敏感性分析}
\end{table}

结果表明，在所选三个 2 月日期的局部工况下，模型对常规幅度原水 NTU 扰动响应较弱；当原水 NTU 增加 50\% 时，部分预测步长出现明显抬升，说明模型捕捉到了较大原水水质突变对未来出厂水浊度的风险传导。矾投加量在这些日期上的局部扰动影响接近 0，可能原因是当前样本处于 XGBoost 树模型的同一分裂区间内，或者矾投加调整对出厂 NTU 的作用主要通过滤后水状态和更长过程滞后体现。

\section{结论}

本文建立了结合质量守恒思想和数据驱动方法的混合动态模型。RTD 层用于描述清水池中历史水体对当前出厂水状态的混合贡献，XGBoost 残差层用于学习混凝、过滤、加药和运行调控中的非线性影响，自回归递推机制用于解决 2026 年 2 月 NTU 缺失导致的状态断裂问题。

验证结果显示，模型在 2 小时预测上具有较高精度，随着预测步长延长，误差整体增大但仍保持可用的动态预测能力。预测结果均低于 1 NTU，说明所选日期的出厂水浊度预测处于达标范围内。敏感性分析显示，较大幅度原水浊度突变可能引起未来 NTU 上升，而矾投加调整在当前局部样本下的直接影响较弱，需要结合滤后水变化和更长滞后进一步分析。

\end{document}
